% Lean: EconLean/Foundations/Correspondence.lean, EconLean/Foundations/Maximum.lean
\chapter{Correspondences and Berge's Maximum Theorem}
\label{chap:foundations}

Sources: Aliprantis--Border (AB), \emph{Infinite Dimensional Analysis}, \S17;
Stokey--Lucas (SL) \S3.3; MWG Appendix M.H, M.K.

A \emph{correspondence} $\Phi : X \twoheadrightarrow Y$ assigns to each $x \in X$ a set
$\Phi(x) \subseteq Y$; in Lean it is a function \texttt{X → Set Y}. Throughout, $X$ and $Y$
are topological spaces and values may be empty unless a hypothesis says otherwise.
Mathlib supplies the definitions of upper and lower hemicontinuity (UHC, LHC); we restate
them in textbook form and build the rest.

\section{Hemicontinuity and graphs}
\label{sec:hemicontinuity}

\begin{definition}[Upper hemicontinuity, textbook form; AB Def.~17.2]
\label{def:uhc}
\lean{Econ.upperHemicontinuousAt_iff_textbook}
\leanok
$\Phi$ is \emph{upper hemicontinuous at $x$} iff for every open $U \supseteq \Phi(x)$ there is a
neighbourhood $N$ of $x$ with $\Phi(x') \subseteq U$ for all $x' \in N$.
\end{definition}
\begin{proof}
\leanok
Mathlib phrases UHC with the neighbourhood filter of the set $\Phi(x)$; the equivalence with the
open-set form is the Mathlib lemma \texttt{upperHemicontinuousAt\_iff\_forall\_isOpen}.
\end{proof}

\begin{definition}[Lower hemicontinuity, textbook form; AB Def.~17.2]
\label{def:lhc}
\lean{Econ.lowerHemicontinuousAt_iff_textbook}
\leanok
$\Phi$ is \emph{lower hemicontinuous at $x$} iff for every open $U$ with
$\Phi(x) \cap U \neq \emptyset$ there is a neighbourhood $N$ of $x$ with
$\Phi(x') \cap U \neq \emptyset$ for all $x' \in N$.
\end{definition}
\begin{proof}
\leanok
This is Mathlib's \texttt{lowerHemicontinuousAt\_iff}.
\end{proof}

\begin{definition}[Continuous correspondence; graph]
\label{def:contCorr}
\lean{Econ.ContinuousCorrespondence, Econ.graph, Econ.mem_graph}
\leanok
\uses{def:uhc, def:lhc}
$\Phi$ is \emph{continuous} if it is UHC and LHC at every point. Its \emph{graph} is
$\operatorname{Gr}\Phi = \{(x, y) \in X \times Y \mid y \in \Phi(x)\}$.
\end{definition}

\begin{lemma}[UHC is local]
\label{lem:uhc_congr}
\lean{Econ.UpperHemicontinuousAt.congr'}
\leanok
\uses{def:uhc}
If $\Phi$ is UHC at $x$ and $\Phi(x') = \Psi(x')$ for all $x'$ in a neighbourhood of $x$, then
$\Psi$ is UHC at $x$.
\end{lemma}
\begin{proof}
\leanok
Use the textbook form. Let $U \supseteq \Psi(x)$ be open. Since $\Phi(x) = \Psi(x)$ (the equality holds
at $x$ itself), $U \supseteq \Phi(x)$, so UHC of $\Phi$ gives a neighbourhood on which
$\Phi(x') \subseteq U$. On the intersection with the neighbourhood where $\Phi = \Psi$, we get
$\Psi(x') \subseteq U$.
\end{proof}

\begin{lemma}[Closed graph gives closed values]
\label{lem:closed_values}
\lean{Econ.isClosed_of_isClosed_graph}
\leanok
\uses{def:contCorr}
If $\operatorname{Gr}\Phi$ is closed then every $\Phi(x)$ is closed.
\end{lemma}
\begin{proof}
\leanok
$\Phi(x)$ is the preimage of $\operatorname{Gr}\Phi$ under the continuous map $y \mapsto (x, y)$.
\end{proof}

\begin{theorem}[UHC with closed values has closed graph; AB Thm~17.10]
\label{thm:closed_graph_of_uhc}
\lean{Econ.isClosed_graph_of_upperHemicontinuous}
\leanok
\uses{def:uhc, def:contCorr}
If $Y$ is regular, $\Phi$ is UHC everywhere and every $\Phi(x)$ is closed, then
$\operatorname{Gr}\Phi$ is closed.
\end{theorem}
\begin{proof}
\leanok
We show the complement is open. Take $(x, y)$ with $y \notin \Phi(x)$. Since $\Phi(x)$ is closed and
$Y$ is regular, the point $y$ and the set $\Phi(x)$ have disjoint neighbourhoods: an open
$W' \supseteq \Phi(x)$ and a neighbourhood $V$ of $y$ with $W' \cap V = \emptyset$. By UHC there is a
neighbourhood $N$ of $x$ on which $\Phi(x') \subseteq W'$. Then $N \times V$ is a neighbourhood of
$(x, y)$, and any $(x', y') \in N \times V$ with $y' \in \Phi(x')$ would put $y'$ in both $W'$ and $V$,
which is impossible. So $N \times V$ misses the graph.
\end{proof}

\begin{theorem}[Intersection with a closed-graph correspondence; AB Thm~17.25]
\label{thm:uhc_inter}
\lean{Econ.UpperHemicontinuousAt.inter_of_isClosed_graph}
\leanok
\uses{def:uhc, def:contCorr}
If $\Phi$ is UHC at $x$ with $\Phi(x)$ compact, and $\Psi$ has closed graph, then
$x' \mapsto \Phi(x') \cap \Psi(x')$ is UHC at $x$.
\end{theorem}
\begin{proof}
\leanok
Let $U$ be open with $\Phi(x) \cap \Psi(x) \subseteq U$. The set $C = \Phi(x) \setminus U$ is compact
(closed subset of a compact set), and $\{x\} \times C$ lies outside $\operatorname{Gr}\Psi$: a point
$y \in C$ in $\Psi(x)$ would lie in $\Phi(x) \cap \Psi(x) \subseteq U$. The complement of
$\operatorname{Gr}\Psi$ is open, so the generalized tube lemma gives open $u \ni x$ and
$v \supseteq C$ with $(u \times v) \cap \operatorname{Gr}\Psi = \emptyset$. Now
$\Phi(x) \subseteq U \cup v$, which is open, so by UHC $\Phi(x') \subseteq U \cup v$ for $x'$ near $x$;
also $x' \in u$ near $x$. For such $x'$ and $y \in \Phi(x') \cap \Psi(x')$: either $y \in U$, as
wanted, or $y \in v$, in which case $(x', y) \in u \times v$ is in the graph of $\Psi$, a contradiction.
\end{proof}

\begin{theorem}[Closed graph and locally bounded values give UHC; AB Thm~17.11]
\label{thm:uhc_of_closed_graph}
\lean{Econ.upperHemicontinuousAt_of_isClosed_graph, Econ.upperHemicontinuous_of_isClosed_graph}
\leanok
\uses{def:uhc, def:contCorr}
If $\operatorname{Gr}\Phi$ is closed and there is a compact $K$ with $\Phi(x') \subseteq K$ for all
$x'$ near $x$, then $\Phi$ is UHC at $x$. In particular, if $Y$ is compact, a closed graph implies UHC.
\end{theorem}
\begin{proof}
\uses{thm:uhc_inter, lem:uhc_congr}
\leanok
The constant correspondence $x' \mapsto K$ is UHC with compact value, so by
Theorem~\ref{thm:uhc_inter} the correspondence $x' \mapsto K \cap \Phi(x')$ is UHC at $x$. Near $x$
it equals $\Phi$ (since $\Phi(x') \subseteq K$), so $\Phi$ is UHC at $x$ by Lemma~\ref{lem:uhc_congr}.
For compact $Y$ take $K = Y$.
\end{proof}

\begin{theorem}[Compact images; AB Lemma~17.8]
\label{thm:compact_image}
\lean{Econ.IsCompact.biUnion_of_upperHemicontinuous}
\leanok
\uses{def:uhc}
If $K \subseteq X$ is compact and $\Phi$ is UHC with compact values at every point of $K$, then
$\bigcup_{x \in K}\Phi(x)$ is compact.
\end{theorem}
\begin{proof}
\leanok
Use the finite-subcover criterion. Let $\{U_i\}$ be an open cover of the union. For each $x \in K$, the
compact set $\Phi(x)$ is covered by finitely many $U_i$, $i \in t_x$. By UHC, $\Phi(x') \subseteq
\bigcup_{i \in t_x} U_i$ for all $x'$ in a neighbourhood $N_x$ of $x$. The $N_x$ cover $K$, so by
compactness finitely many $N_{x_1}, \ldots, N_{x_m}$ do. Then $\bigcup_k t_{x_k}$ is a finite index set
whose $U_i$ cover every $\Phi(x')$, $x' \in K$.
\end{proof}

\begin{proposition}[Closed graph, sequential form; AB 17.16]
\label{prop:closed_graph_seq}
\lean{Econ.isClosed_graph_iff_seq}
\leanok
\uses{def:contCorr}
If $X$ and $Y$ are first countable, $\operatorname{Gr}\Phi$ is closed iff whenever $x_n \to x$,
$y_n \to y$ and $y_n \in \Phi(x_n)$ for all $n$, we have $y \in \Phi(x)$.
\end{proposition}
\begin{proof}
\leanok
($\Rightarrow$) $(x_n, y_n) \to (x, y)$ and every term lies in the closed graph, so the limit does.
($\Leftarrow$) In a first-countable space closed equals sequentially closed. A sequence $q_n$ in the
graph converging to $p$ has coordinates $q_{n,1} \to p_1$ and $q_{n,2} \to p_2$ with
$q_{n,2} \in \Phi(q_{n,1})$; the hypothesis gives $p \in \operatorname{Gr}\Phi$.
\end{proof}

\begin{proposition}[Convergent subsequences under UHC; AB 17.20]
\label{prop:uhc_subseq}
\lean{Econ.UpperHemicontinuousAt.exists_tendsto_subseq}
\leanok
\uses{def:uhc}
Let $Y$ be a (pseudo)metric space, $\Phi$ UHC at $x$ with $\Phi(x)$ compact. If $x_n \to x$ and
$y_n \in \Phi(x_n)$, then some subsequence $y_{\varphi(n)}$ converges to a point of $\Phi(x)$.
\end{proposition}
\begin{proof}
\leanok
First, $\Phi(x) \neq \emptyset$: otherwise UHC with $U = \emptyset$ forces $\Phi(x_n) = \emptyset$ for
large $n$, contradicting $y_n \in \Phi(x_n)$. Next, $d(y_n, \Phi(x)) \to 0$: for $\varepsilon > 0$ the set
$\{y \mid d(y, \Phi(x)) < \varepsilon\}$ is open and contains $\Phi(x)$, so by UHC it contains
$\Phi(x_n) \ni y_n$ for large $n$. Since $\Phi(x)$ is compact and nonempty, choose nearest points
$z_n \in \Phi(x)$ with $d(y_n, z_n) = d(y_n, \Phi(x))$. By compactness a subsequence
$z_{\varphi(n)} \to y \in \Phi(x)$. Finally
$d(y_{\varphi(n)}, y) \le d(y_{\varphi(n)}, z_{\varphi(n)}) + d(z_{\varphi(n)}, y) \to 0$.
\end{proof}

\section{Berge's maximum theorem}
\label{sec:berge}

Let $f : X \times Y \to \R$ (written $f(x, y)$) and $\Phi : X \twoheadrightarrow Y$.

\begin{definition}[Value function and argmax]
\label{def:maxValue}
\lean{Econ.maxValue, Econ.argmax}
\leanok
The \emph{value function} is $m(x) = \sup\{f(x, y) \mid y \in \Phi(x)\}$ and the \emph{argmax} is
$\mu(x) = \{y \in \Phi(x) \mid \forall y' \in \Phi(x),\ f(x, y') \le f(x, y)\}$.
In Lean the supremum of an empty or unbounded set of reals is $0$, so every result below assumes
nonempty compact values, under which $m(x)$ is a genuine maximum.
\end{definition}

\begin{lemma}[The maximum is attained]
\label{lem:isGreatest_maxValue}
\lean{Econ.continuous_section, Econ.isGreatest_maxValue}
\leanok
\uses{def:maxValue}
If $f$ is jointly continuous then each section $f(x, \cdot)$ is continuous, and if $\Phi(x)$ is nonempty
and compact then $m(x)$ is the greatest element of $f(x, \Phi(x))$.
\end{lemma}
\begin{proof}
\leanok
The section is $f$ composed with $y \mapsto (x, y)$. The image of the nonempty compact $\Phi(x)$
under the continuous $f(x, \cdot)$ is nonempty and compact in $\R$, so contains its supremum.
\end{proof}

\begin{lemma}[Argmax as a superlevel set]
\label{lem:argmax_eq}
\lean{Econ.argmax_eq}
\leanok
\uses{lem:isGreatest_maxValue}
Under the same hypotheses, $\mu(x) = \Phi(x) \cap \{y \mid m(x) \le f(x, y)\}$.
\end{lemma}
\begin{proof}
\leanok
If $y \in \mu(x)$, pick $y_0 \in \Phi(x)$ attaining $m(x)$ (Lemma~\ref{lem:isGreatest_maxValue});
then $m(x) = f(x, y_0) \le f(x, y)$. Conversely if $m(x) \le f(x, y)$ then for every feasible $y'$,
$f(x, y') \le m(x) \le f(x, y)$.
\end{proof}

\begin{theorem}[Berge, upper half]
\label{thm:usc_maxValue}
\lean{Econ.upperSemicontinuous_maxValue}
\leanok
\uses{lem:isGreatest_maxValue, def:uhc}
If $f$ is continuous and $\Phi$ is UHC with nonempty compact values, then $m$ is upper semicontinuous.
\end{theorem}
\begin{proof}
\leanok
Fix $x$ and $c > m(x)$. Every $y \in \Phi(x)$ has $f(x, y) \le m(x) < c$, so $\{x\} \times \Phi(x)$ lies in
the open set $W = \{(x', y) \mid f(x', y) < c\}$. By the tube lemma (with $\{x\}$ and $\Phi(x)$ compact)
there are open $u \ni x$ and $v \supseteq \Phi(x)$ with $u \times v \subseteq W$. By UHC,
$\Phi(x') \subseteq v$ for $x'$ near $x$, and also $x' \in u$. For such $x'$ take a maximizer $y'$
with $f(x', y') = m(x')$; then $(x', y') \in u \times v \subseteq W$, so $m(x') < c$.
\end{proof}

\begin{theorem}[Berge, lower half]
\label{thm:lsc_maxValue}
\lean{Econ.lowerSemicontinuous_maxValue}
\leanok
\uses{lem:isGreatest_maxValue, def:lhc}
If $f$ is continuous and $\Phi$ is LHC with nonempty compact values, then $m$ is lower semicontinuous.
\end{theorem}
\begin{proof}
\leanok
Fix $x$ and $c < m(x)$. Pick a maximizer $y \in \Phi(x)$ with $f(x, y) = m(x) > c$. The set
$\{(x', y') \mid c < f(x', y')\}$ is open and contains $(x, y)$, so it contains a box $u \times v$ with
$x \in u$, $y \in v$. Since $v$ meets $\Phi(x)$ at $y$, LHC gives, for $x'$ near $x$ (and in $u$), a
feasible $y' \in \Phi(x') \cap v$. Then $c < f(x', y') \le m(x')$.
\end{proof}

\begin{theorem}[Berge: continuity of the value]
\label{thm:continuous_maxValue}
\lean{Econ.continuous_maxValue}
\leanok
\uses{thm:usc_maxValue, thm:lsc_maxValue, def:contCorr}
If $f$ is continuous and $\Phi$ is a continuous correspondence with nonempty compact values, then $m$ is
continuous.
\end{theorem}
\begin{proof}
\leanok
A real function is continuous iff it is both lower and upper semicontinuous; apply
Theorems~\ref{thm:lsc_maxValue} and~\ref{thm:usc_maxValue}.
\end{proof}

\begin{lemma}[Argmax nonempty and compact]
\label{lem:argmax_props}
\lean{Econ.argmax_nonempty, Econ.isCompact_argmax}
\leanok
\uses{lem:argmax_eq}
If $f$ is continuous and $\Phi(x)$ is nonempty and compact, then $\mu(x)$ is nonempty and compact.
\end{lemma}
\begin{proof}
\leanok
Nonempty: a continuous function attains its maximum on a nonempty compact set (Weierstrass).
Compact: by Lemma~\ref{lem:argmax_eq}, $\mu(x)$ is the intersection of the compact $\Phi(x)$ with the
closed set $\{y \mid m(x) \le f(x, y)\}$.
\end{proof}

\begin{theorem}[Berge: UHC of the argmax]
\label{thm:uhc_argmax}
\lean{Econ.upperHemicontinuous_argmax}
\leanok
\uses{thm:continuous_maxValue, thm:uhc_inter, lem:argmax_eq, lem:uhc_congr}
If $f$ is continuous and $\Phi$ is a continuous correspondence with nonempty compact values, then $\mu$
is UHC.
\end{theorem}
\begin{proof}
\leanok
Let $\Psi(x) = \{y \mid m(x) \le f(x, y)\}$. Since $m$ is continuous
(Theorem~\ref{thm:continuous_maxValue}) and $f$ is continuous, $\operatorname{Gr}\Psi =
\{(x, y) \mid m(x) \le f(x, y)\}$ is closed. By Theorem~\ref{thm:uhc_inter}, $\Phi \cap \Psi$ is UHC at
each $x$, and by Lemma~\ref{lem:argmax_eq} $\Phi \cap \Psi = \mu$ everywhere, so $\mu$ is UHC
(Lemma~\ref{lem:uhc_congr}).
\end{proof}

\begin{theorem}[Berge's maximum theorem; AB 17.31, SL 3.6, MWG M.K.6]
\label{thm:berge}
\lean{Econ.berge_maximum_theorem}
\leanok
\uses{lem:isGreatest_maxValue, thm:continuous_maxValue, lem:argmax_props, thm:uhc_argmax}
Let $f : X \times Y \to \R$ be continuous and $\Phi$ a continuous correspondence with nonempty compact
values. Then (1) $m(x)$ is attained for every $x$; (2) $m$ is continuous; (3) $\mu$ has nonempty compact
values; (4) $\mu$ is UHC.
\end{theorem}
\begin{proof}
\leanok
Collect Lemma~\ref{lem:isGreatest_maxValue}, Theorem~\ref{thm:continuous_maxValue},
Lemma~\ref{lem:argmax_props} and Theorem~\ref{thm:uhc_argmax}.
\end{proof}

\begin{corollary}[Unique maximizers vary continuously]
\label{cor:continuous_argmax_singleton}
\lean{Econ.continuous_of_argmax_eq_singleton}
\leanok
\uses{thm:uhc_argmax}
In the setting of Theorem~\ref{thm:berge}, if $\mu(x) = \{g(x)\}$ for every $x$, then $g$ is continuous.
\end{corollary}
\begin{proof}
\leanok
A singleton-valued correspondence $x \mapsto \{g(x)\}$ is UHC iff $g$ is continuous (Mathlib's
\texttt{upperHemicontinuous\_singleton\_iff}); apply Theorem~\ref{thm:uhc_argmax}.
\end{proof}

\begin{remark}[LHC cannot be dropped]
\label{rem:jumpCorr}
\lean{Econ.jumpCorr}
\leanok
\uses{thm:continuous_maxValue}
The Lean file also checks a counterexample: $\Phi(0) = [0, 1]$, $\Phi(x) = \{0\}$ for $x \neq 0$, with
$f(x, y) = y$. $\Phi$ has closed graph $(\{0\} \times [0,1]) \cup (\R \times \{0\})$ with values inside
$[0, 1]$, so it is UHC with nonempty compact values (Theorem~\ref{thm:uhc_of_closed_graph}). It is not
LHC at $0$: the open set $(1/2, \infty)$ meets $\Phi(0)$ but misses $\Phi(x)$ for $x \neq 0$. And
$m(0) = 1$ while $m(x) = 0$ for $x \neq 0$, so $m(1/(n+1)) = 0 \not\to 1$. A second check verifies
all hypotheses of Theorem~\ref{thm:berge} for $\max_{y \in [-1,1]} -(y - x)^2$, so the theorem is not
vacuous.
\end{remark}
